\section{Nested bubble insertions}
\label{sec:nested-bubble}

This chapter covers L2E2G2 and L2E3G4, L2E3G6, L2E3G7, with the fixed-leg
maps of Section~\ref{sec:scheme}. The internal bubble is subtracted through
second order in its external momentum. Integrating that already subtracted
subgraph does not change its finite Taylor conditions. In particular, its
finite kinetic Taylor term must be retained.

\subsection{The finite one-loop bubble}

Choose the inserted-edge momentum $k$ without physical external-momentum
dependence, and let $q$ be the bubble loop variable. Put $r=k^2$. Then
\begin{align*}
 B_R(r)&=\int_q(1-t_k^2)[D(q)D(q+k)]\\
 &=\int_0^1\left\{(M+hr)\log\left(1+\frac{hr}{M}\right)-hr\right\}\,dx,
 \qquad h=x(1-x).
\end{align*}
This formula follows with a common two-edge Schwinger cutoff: write the
parameters as $\rho(1-x),\rho x$. Their Jacobian is $\rho$ and the
Gaussian factor is $\rho^{-3}$. The Taylor-subtracted radial kernel is
\[
 \int_{T(x)}^\infty\frac{d\rho}{\rho^2}e^{-M\rho}
       (e^{-hr\rho}-1+hr\rho),\qquad
 T(x)=\frac{\eta}{\min(x,1-x)}.
\]
It is nonnegative and is controlled by its zero-cutoff value
$K(hr)$, where
\[
 K(z)=(M+z)\log(1+z/M)-z.
\]
The parameter regulator limit is therefore legitimate. Independent momentum
absolute convergence and its relation to this formula are proved below.

The finite conditions and second derivative are
\[
 B_R(0)=B_R'(0)=0,\qquad
 B_R''(r)=\int_0^1\frac{h^2}{M+hr}\,dx.
\]
These conditions uniquely fix $B_R$. In particular the term
$-r\int_0^1h\,dx=-r/6$ is part of the original finite BPH subtraction;
discarding it would change the two-loop answer. An equivalent closed form is
\[
 B_R(r)=\frac{(r+4M)^{3/2}}{3\sqrt r}
            \operatorname{arsinh}\sqrt{\frac r{4M}}
                -\frac{4M}{3}-\frac{4r}{9},
\]
with the continuous value at zero. For $|r|<4M$ the convergent expansion is
\[
 \frac{B_R(r)}M=
 \sum_{n\ge2}\frac{(-1)^n(r/M)^n(n!)^2}
                         {n(n-1)(2n+1)!},
\]
whose first term is $(r/M)^2/60$. Since
$0\le K(z)\le z\log(1+z/M)$, $h\le1/4$, and $\int_0^1h\,dx=1/6$,
\[
 0\le B_R(r)\le\frac r6\log\left(1+\frac r{4M}\right)
                    \le\frac r6\log\left(1+\frac rM\right).
\]

\subsection{A uniform original-momentum subgraph norm}

Write $K=|k|$. The explicit bound required for Fubini at two loops is
\[
 \int_q\left|(1-t_k^2)[D(q)D(q+k)]\right|
 \le K^2\left\{218+5\log\left(1+\frac{4K}{m}\right)\right\}.
\]
We give its derivation, rather than relying only on power counting.
For $K\le m$, the third-order Taylor integral remainder, the derivative
constant $c_3=72$, and
$\langle q+tk\rangle\ge\langle q\rangle/2$ imply the pointwise bound
\[
 \left|(1-t_k^2)[D(q)D(q+k)]\right|
       \le384K^3\langle q\rangle^{-7}.
\]
The factor $384$ is $(72/6)2^5$. Since
\[
 \int_q\langle q\rangle^{-7}=\frac8{15m},
\]
the norm is at most $(1024/5)K^3/m\le(1024/5)K^2<218K^2$.

For $K\ge m$, split at $R_*=2(K+m)$. On $|q|\ge R_*$ the same remainder
bound integrates to
$384K^3/R_*\le192K^2$. In the ball $|q|\le R_*$, use the following
separate integrable bounds, which are allowed on this finite region:
\begin{align*}
 \int_{|q|\le R_*}D(q)D(q+k)
 &\le\frac{R_*(R_*+K)}2\le10K^2,\\
 \int_{|q|\le R_*}D(q)^2&\le\frac{R_*^2}{2}\le8K^2,\\
 \int_{|q|\le R_*}|D(q)\,k\cdot\nabla D(q)|
 &\le2KR_*\le8K^2.
\end{align*}
The first estimate is Cauchy--Schwarz, with the shifted ball enclosed in
radius $R_*+K$. The quadratic Taylor term is bounded by
\begin{align*}
 5K^2\int_0^{R_*}\frac{Q^5\,dQ}{(Q^2+m^2)^3}
 &\le\frac52K^2\log\left(1+\frac{16K^2}{m^2}\right)\\
 &\le5K^2\log\left(1+\frac{4K}{m}\right).
\end{align*}
Here $\operatorname{area}(S^5)=\pi^3$, so the normalized radial measure
really is $Q^5dQ$. Adding $192+10+8+8$ proves the constant $218$.

\subsection{Original twelve-dimensional convergence and the regulator bridge}

At fixed $k$, the subgraph remainder is independent of physical external
momenta. The whole-graph Taylor operator therefore acts only on the shifted
outer propagators. For the two-point insertion the remaining factor is
\[
 E_2(k,p)=D(k)^2(1-t_p^2)D(k+p).
\]
If $K\ge2|p|$, the third-order remainder gives
\[
 |E_2(k,p)|\le384|p|^3\langle k\rangle^{-9}.
\]
For the three-point insertion, with adjacent momenta $P,Q$, the outer
factor is
\[
 E_3(k;P,Q)=D(k)^2\{D(k+P)D(k+Q)-D(k)^2\}.
\]
If $K\ge2\max(|P|,|Q|)$, a first-order physical-ray remainder gives
\[
 |E_3(k;P,Q)|\le64(|P|+|Q|)\langle k\rangle^{-9}.
\]
Multiplying either bound by the subgraph $q$-norm just proved, and using
the outer radial measure $K^5dK$, bounds the tail by a constant times
$K^{-2}\log K$. It is integrable. On every finite $K$ interval, positive
mass bounds the denominators and the displayed subgraph norm is finite.
Thus the complete original twelve-dimensional recursive expression is
absolutely integrable.

All these controls also hold for $D_\eta$, including the Taylor remainder
and the finite-ball bounds. The positive Schwinger moment estimates in
Section~\ref{sec:scheme} give the same constants uniformly in $\eta$.
At positive cutoff, Fubini permits the $q$ integration first. Its limit is
$B_R(k^2)$ by the one-loop parameter formula above; the two-loop majorant
permits the same limit under the remaining integral. Equivalently,
momentum dominated convergence gives the original ordinary $RI$ integral,
and the subgraph parameter limit gives the reduced formula. They coincide.
No separate unregulated divergent bare or counterterm integral is used.
The argument is uniform on bounded external momenta and $m\ge m_0>0$,
by using $m_0$ in every majorant from the outset.

This factorization is routing-sensitive. If one changes to
$k'=k+\alpha p$, the overall Taylor derivatives must also act on
$B_R((k'-\alpha p)^2)$ and on every other resulting external-momentum
dependence. Substituting $B_R(k'^2)$ instead loses finite local Taylor
terms. Similarly a subgraph loop shift applies to the complete subtracted
combination. The catalog routings can be put in the present form by
external-momentum-independent loop changes with absolute determinant one.

\subsection{Two-point insertion: positive one-dimensional integral}

Put $s=p^2$, $r=k^2$, and
\[
 R=r+M,\qquad A=R+s,\qquad W=\sqrt{(R-s)^2+4Ms}.
\]
Normalized six-dimensional angular averaging uses the density
$\frac8{3\pi}(1-z^2)^{3/2}dz$ on $-1\le z\le1$. Direct integration gives
\[
 \langle D(k+p)\rangle_{S^5}
        =\frac{4(A+2W)}{3(A+W)^2}.
\]
The odd Taylor term averages to zero, and
$\langle(k\cdot p)^2\rangle=rs/6$. Define
\begin{align*}
 F_2(r,s)&=\langle D(k+p)\rangle_{S^5}-\frac1R
                    -s\left(-\frac1{R^2}+\frac{2r}{3R^3}\right),\\
 J_{\mathrm{b},2}(s;M)
       &=\frac12\int_0^\infty\frac{r^2B_R(r)}{R^2}F_2(r,s)\,dr.
\end{align*}
The radial Jacobian is $K^5dK=\frac12r^2dr$.

For a manifestly stable sign, set $\mu=M/R$, $\sigma=s/R$, and
\[
 d=\frac{\sigma-1+\sqrt{(1-\sigma)^2+4\mu\sigma}}2,
 \qquad 0\le d\le\sigma.
\]
Algebra using $d^2+(1-\sigma)d-\mu\sigma=0$ yields
\[
 F_2(r,s)=\frac{d^2\{(1+2\mu)d^2+6\mu d+3\mu\}}
                     {3R(1+d)^2(d+\mu)}\ge0.
\]
Moreover,
\begin{align*}
 &3(1+d)^2(d+\mu)-\{(1+2\mu)d^2+6\mu d+3\mu\}\\
 &\hspace{25mm}=3d^3+(5+\mu)d^2+3d\ge0.
\end{align*}
Consequently $F_2\le d^2/R\le s^2/R^3$. The exact radial constant is
\[
 \int_0^\infty\frac{z^3\log(1+z)}{(1+z)^5}\,dz=\frac{25}{48}.
\]
For example, it is the logarithmic derivative of the beta integral
$B(4,1)=1/4$. Combining it with the bound on $B_R$ proves
\[
 0\le J_{\mathrm{b},2}(s;M)\le\frac{25s^2}{576M}.
\]
The mass dimension is two:
$J_{\mathrm{b},2}(s;M)=M J_{\mathrm{b},2}(s/M;1)$.
It vanishes with its first $s$ derivative at zero, as required by the
overall second-order momentum subtraction.

A stronger large-scale estimate uses
\[
 \frac1R-\langle D(k+p)\rangle_{S^5}
   =\frac{\sigma+2d+3d^2}{3R(1+d)^2}\ge0,
 \qquad F_2\le\frac{s}{R^2}.
\]
Split the radial integral at $r=s$. On the lower part use $F_2\le s/R^2$,
and on the upper part use $F_2\le s^2/R^3$. With
$r^3/R^4\le1/R$, $r^3/R^5\le1/R^2$, and
$L=\log(1+s/M)$, direct integration gives
\[
 J_{\mathrm{b},2}(s;M)\le
 \min\left\{\frac{25s^2}{576M},
       \frac{sL^2}{24}+\frac{s^2(L+1)}{12(M+s)}\right\}.
\]
This is at most of order $s\log^2s$ at a large common external scale.

\subsection{Three-point insertion: exact angular reduction and its bound}

Let $a=P^2$, $b=Q^2$, $c=(P-Q)^2$; $a,b$ are adjacent to the inserted
edge and $c$ is opposite it. The three fixed-label graphs use
$(P,Q)=(p_1,-p_2),(p_1,-p_3),(p_2,-p_3)$, respectively. Their physical
domain is $a,b,c\ge0$ and $ab-(a+b-c)^2/4\ge0$.
The ordinary finite integral is
\[
 J_{\mathrm{b},3}(a,b,c;M)
 =\int_k D(k)^2\{D(k+P)D(k+Q)-D(k)^2\}B_R(k^2).
\]
Use a Feynman parameter $u\in[0,1]$ and put
\[
 S=ua+(1-u)b,\qquad T=u(1-u)c.
\]
The parameter Jacobian is one and
\[
 D(k+P)D(k+Q)
 =\int_0^1\frac{du}{\{(k+uP+(1-u)Q)^2+M+T\}^2}.
\]
For $r=k^2$, let $R=r+M$, $A=R+S$, and
\[
 W=\sqrt{(R-S)^2+4MS+4rT}.
\]
The exact angular identity
\[
 \left\langle\frac1{(A+B\cos\theta)^2}\right\rangle_{S^5}
   =\frac4{(A+\sqrt{A^2-B^2})^2},\qquad
 B^2=4r(S-T),
\]
can also be obtained by differentiating the preceding first-power angular
integral at fixed $B$. It gives the two-dimensional finite representation
\[
 J_{\mathrm{b},3}
 =\frac12\int_0^\infty\frac{r^2B_R(r)}{R^2}
       \int_0^1\left\{\frac4{(A+W)^2}-\frac1{R^2}\right\}\,du\,dr.
\]
The Gram condition gives $S-T=|uP+(1-u)Q|^2\ge0$, and $0\le T\le S$.
Thus $2R\le A+W\le2(R+S)$, whence
\[
 0\le\frac1{R^2}-\frac4{(A+W)^2}\le\frac{2S}{R^3}.
\]
Using the same radial constant $25/48$ and
$\int_0^1S\,du=(a+b)/2$, we obtain
\[
 -\frac{25(a+b)}{576M}\le J_{\mathrm{b},3}(a,b,c;M)\le0.
\]
This kernel has mass dimension zero. Adjacent exchange $a\leftrightarrow b$
is a symmetry, whereas moving the opposite leg generally changes the value.

The nonnegative angular average also bounds the difference by $1/R^2$.
Let $H=\max(a,b)$, $L=\log(1+2H/M)$. Split at $r=2H$, use $1/R^2$
below and $2H/R^3$ above, and integrate as in the two-point case. The
result is
\[
 |J_{\mathrm{b},3}|\le
 \min\left\{\frac{25(a+b)}{576M},
        \frac{L^2}{24}+\frac{H(L+1)}{6(M+2H)}\right\}.
\]
At $H=0$, the physical Gram domain forces all momenta to vanish and the
continuous bound is zero. At large common scale it grows at most as
$\log^2$. Both insertion bounds are compact-uniform for bounded external
momenta and positive mass lower bound. Their sharp constants refer to the
fixed-sign angularly reduced densities; the original momentum $L^1$ proof
uses the separate, coarser subgraph estimate given above.
